\section{Independence, products, and conditional expectation}\label{sec:independence}

\subsection{Product laws}

\begin{definition}[Finite product law]\label{def:product-law}
Let \(L=(m,w,W)\) and \(M=(n,v,V)\).  Their product law has outcomes
\((i,j)\in\Fin(m)\times\Fin(n)\), weight \(w_iv_j\), and total weight
\[
 \sum_{i<m}\sum_{j<n}w_iv_j=WV.
\]
An observable depending only on the first coordinate is written \(X(i)\), and
one depending only on the second is written \(Y(j)\).
\end{definition}

\begin{theorem}[Multiplicativity under independence]\label{thm:ind-product}
On the product law,
\[
 \E[X(i)Y(j)]=\E_L[X]\E_M[Y].
\]
More generally, for finite lists of independent coordinates and rational
functions \(g_r\) of one coordinate,
\[
 \E\left[\prod_{r<k}g_r(X_r)\right]
 =\prod_{r<k}\E[g_r(X_r)].
\]
\end{theorem}

\begin{proof}
For two coordinates,
\begin{align*}
 WV\E[XY]
 &=\sum_{i<m}\sum_{j<n}w_iv_jX_iY_j\\
 &=\left(\sum_{i<m}w_iX_i\right)
   \left(\sum_{j<n}v_jY_j\right),
\end{align*}
by finite distributivity and Fubini.  Divide by \(WV>0\).  The multi-coordinate
formula follows by quantifier-free induction on \(k\), using a mixed-radix code
for product outcomes and the two-coordinate identity at each successor step.
\end{proof}

\begin{corollary}[Variance of an independent sum]\label{cor:ind-var-sum}
If \(X,Y\) are independent, then \(\Cov(X,Y)=0\) and
\[
 \Var(X+Y)=\Var(X)+\Var(Y).
\]
For a finite independent family,
\[
 \Var\left(\sum_{j<r}X_j\right)=\sum_{j<r}\Var(X_j).
\]
\end{corollary}

\begin{proof}
Multiplicativity gives \(\E[XY]=\E[X]\E[Y]\), so covariance is zero.  Insert
this into the variance expansion.  Induct on the number of summands; a new
coordinate is independent of every function of the preceding product
coordinates, a fact verified by another finite product factorization.
\end{proof}

\begin{example}[Expectation is not multiplicative in general]\label{ex:not-multiplicative}
Let \(X=Y\) be uniform on \(\{-1,1\}\).  Then
\(\E[X]=\E[Y]=0\), while \(XY=1\) identically and \(\E[XY]=1\).
Thus multiplicativity is a consequence of independence, not a basic property
of expectation.
\end{example}

\begin{example}[Uncorrelated does not imply independent]\label{ex:uncorrelated}
Let \(X\) be uniform on \(\{-1,0,1\}\) and let \(Y=X^2\).  Then
\[
 \E[X]=0,
 \qquad \E[XY]=\E[X^3]=0,
 \qquad \Cov(X,Y)=0.
\]
But \(Y\) is a deterministic function of \(X\), and the event \(Y=0\) is
exactly the event \(X=0\); the joint law does not factor.  All values are
finite rationals, so this counterexample is internal to \(\ipra\).
\end{example}

\subsection{Finite partitions and conditional expectation}

\begin{definition}[Finite conditioning partition]\label{def:partition}
A finite partition is a primitive-recursive map
\(c:\Fin(m)\to\Fin(r)\).  Its cell masses and weighted \(X\)-sums are
\[
 M_j=\sum_{i<m}w_i\ind_{c(i)=j},\qquad
 S_j=\sum_{i<m}w_iX_i\ind_{c(i)=j}.
\]
A conditional-expectation code is a cellwise rational value
\[
 h_j=\begin{cases}S_j/M_j,&M_j>0,\\0,&M_j=0.\end{cases}
\]
The observable \(\E[X\mid c]\) takes value \(h_{c(i)}\) at outcome \(i\).
\end{definition}

The arbitrary value on a zero-mass cell is harmless: its total weight is zero.
The branch \(M_j=0\) is decidable, and division is only performed in the
positive branch.

\begin{theorem}[Cell integral characterization]\label{thm:cell-characterization}
For every cell \(j<r\),
\[
 \sum_{i<m}w_i\E[X\mid c](i)\ind_{c(i)=j}
 =\sum_{i<m}w_iX_i\ind_{c(i)=j}.
\]
Equivalently,
\[
 \E[\E[X\mid c]\ind_{c=j}]=\E[X\ind_{c=j}].
\]
\end{theorem}

\begin{proof}
If \(M_j=0\), zero-sum elimination shows every weight in the cell is zero, so
both sides vanish.  If \(M_j>0\), the left numerator is
\(h_jM_j=(S_j/M_j)M_j=S_j\).  The proof is a decidable split on the natural
number \(M_j\).
\end{proof}

\begin{theorem}[Law of total expectation]\label{thm:total-expectation}
For every finite partition,
\[
 \E[\E[X\mid c]]=\E[X].
\]
More generally, if \(g\) is constant on each cell, then
\[
 \E[g\,\E[X\mid c]]=\E[gX].
\]
\end{theorem}

\begin{proof}
Sum the cell characterization over \(j<r\).  For each outcome exactly one cell
indicator is one, so both double sums collapse to the corresponding full sum.
For the generalized identity, multiply the equation for cell \(j\) by the
cell value \(g_j\) before summing.
\end{proof}

\begin{theorem}[Tower property for nested finite partitions]\label{thm:tower}
Suppose \(d:\Fin(m)\to\Fin(s)\) refines
\(c:\Fin(m)\to\Fin(r)\): there is \(q:\Fin(s)\to\Fin(r)\) with
\(c=q\circ d\).  Then
\[
 \E[\E[X\mid d]\mid c]=\E[X\mid c]
\]
almost surely.
\end{theorem}

\begin{proof}
Both sides are constant on each \(c\)-cell.  By
\Cref{thm:cell-characterization,thm:total-expectation}, their weighted integral
on that cell equals the integral of \(X\) on the cell.  If the cell has
positive mass, equality of integral and constancy imply equality of the two
cell values after cancellation.  If the cell has zero mass, disagreement is
confined to a null cell.  The refinement relation makes the regrouping of
\(d\)-cells into a \(c\)-cell a finite Fubini identity.
\end{proof}

\begin{theorem}[Conditional Jensen with certificates]\label{thm:conditional-jensen}
Let \(\phi\) admit a finite Jensen certificate at the conditional mean in each
positive-mass cell.  Then
\[
 \phi(\E[X\mid c])\le \E[\phi(X)\mid c]
\]
almost surely.
\end{theorem}

\begin{proof}
Apply certified Jensen separately to the normalized finite law on each
positive-mass cell.  A zero-mass cell may be ignored.  Combining the finite
list of cellwise inequalities yields the almost-sure statement.
\end{proof}

\subsection{Fubini and effective products}

Finite Fubini is exact.  For countable or continuous product presentations, a
Fubini theorem needs an absolute-integrability modulus.  The constructive
pattern is as follows: truncate both coordinates at stages computed from the
tail modulus; apply finite Fubini; and bound the discarded absolute mass.
This will be formalized in \Cref{sec:integration}.  Independence then implies
multiplicativity whenever the product and the two factors are effectively
integrable.
